\documentclass[11pt]{article}

\usepackage[utf8]{inputenc}
\usepackage[T1]{fontenc}
\usepackage{lmodern}
\usepackage{geometry}
\usepackage{amsmath,amssymb,amsfonts,amsthm,mathtools}
\usepackage{bm}
\usepackage{enumitem}
\usepackage{microtype}
\usepackage{hyperref}
\usepackage{titlesec}
\usepackage{booktabs}
\usepackage{array}
\usepackage{setspace}
\usepackage{mathrsfs}
\usepackage{physics}

\geometry{margin=1in}
\hypersetup{colorlinks=true, linkcolor=black, urlcolor=blue, citecolor=black}
\setlist[enumerate]{topsep=0.4em,itemsep=0.35em}
\setlist[itemize]{topsep=0.3em,itemsep=0.25em}
\titleformat{\section}{\large\bfseries}{\thesection.}{0.5em}{}
\titleformat{\subsection}{\normalsize\bfseries}{\thesubsection.}{0.5em}{}
\setstretch{1.06}

\newcommand{\Aether}{\AE{}ther}
\newcommand{\Aflow}{\AE{}ther-flow}
\newcommand{\TheoryName}{\Aflow{} Interpretation of Relativity}
\newcommand{\TheoryProgram}{\Aether{} / \Aflow{} framework}
\newcommand{\Order}{\mathcal{S}}
\newcommand{\Sub}{\mathcal{A}}
\newcommand{\Ph}{\Phi}
\newcommand{\PhiSrc}{\Phi_{\mathrm{src}}}

\newtheorem{definition}{Definition}
\newtheorem{proposition}{Proposition}
\newtheorem{remark}{Remark}

\title{The \TheoryName{}\\[0.4em]
Ontology: Foundations}
\author{}
\date{}
\makeatletter
\newcommand{\AetherFlowAPATitlePage}{%
  \pagenumbering{arabic}%
  \setcounter{page}{1}%
  \begin{center}
    \vspace*{1.0in}
    {\Large\bfseries \@title\par}
    \vspace{1.2in}
    {\normalsize Project creator: Alexander Samuel Ricciardi\par}
    \vspace{1.2em}
    {\normalsize Model-assisted research and drafting tools: GPT-5.4 and GPT-5.5.\par}
    \vspace{1.2em}
    \begin{minipage}{0.78\textwidth}
      \centering
      {\small Project provenance: this document is part of the current managed ontology package. This title page records project-origin attribution and model-assisted drafting provenance; it does not change the article's scientific content or claim status.\par}
    \end{minipage}
  \end{center}
  \clearpage
}
\makeatother


\begin{document}

\AetherFlowAPATitlePage

\begin{abstract}
This paper provides the foundations statement of \TheoryName{}. At source level, the canonical mainline research architecture is continuum-first and deliberately narrow: it takes a four-dimensional smooth source manifold \(\Sub\) as explicit primitive debt and retains \(\PhiSrc\) only as an unresolved abstract source-order or evolution slot. The source manifold is not identified with physical spacetime, and no domain, algebraic structure, invertibility, stochastic status, gauge status, or physical-time meaning is assigned to \(\PhiSrc\).

Within the broader \TheoryProgram{}, \TheoryName{} is the exact-closure theory developed in the active manuscript sequence. It preserves the \Aether{} / \Aflow{} vocabulary as a proposed interpretation while taking the effective gravitational dynamics to be exactly those of general relativity with ordinary matter coupling. The present manuscript therefore does not claim a completed microphysical derivation of Einsteinian gravity from source variables. Its narrower purpose is to state the selected source architecture, preserve its explicit primitive debt, identify the exact relation to GR, and separate canonical source structure from historical interpretation and still-open physical bridges.

Within the flagship exact-closure package, the companion \emph{Exact-Closure Sequence Overview} is the front door and the \emph{Exact-Closure Note} is the short anchor. The present paper is the first manuscript in the five-paper modular full statement. Its task is to place ontology before derivational ambition, so that exact closure is read first and any later substrate derivation is kept explicitly secondary.
%
\end{abstract}

\section{Introduction}

The \TheoryProgram{} seeks to turn the \Aether{} / \Aflow{} ontology into a mathematically controlled theory. Within that broader program, \TheoryName{} is the exact relativistic theory developed in the active manuscript sequence. In that precise sense, it is the correct entry point for the current paper set.

\paragraph{Framework and claim boundary.}
\TheoryProgram{} denotes the broader \Aether{} / \Aflow{} research interpretation. Its selected mainline source architecture is a four-dimensional smooth manifold \(\Sub\), assumed rather than derived, together with the unresolved slot \(\PhiSrc\). The familiar language of deeper substrate, flow, experiential space, S-time, and observed expansion records motivating interpretation and open bridge questions; it does not add physical spacetime, a slice map, a clock, a metric, or dynamics to the canonical source package. Within that framework, \TheoryName{} denotes the exact-closure benchmark in which the effective gravitational dynamics are adopted to be exactly Einsteinian with universal matter coupling. In the active sequence, \emph{adoption} of the benchmark means use of that established target-side dynamics without claiming source derivation, while \emph{derivation} is reserved for a first-principles recovery from explicit source variables.

The need for such an entry point is structural rather than merely editorial. The project did not begin as a deviation-first program. Its initial impulse was to propose a deeper account of what observed space, temporal ordering, and relativistic phenomenology might be \emph{about}. The present manuscript restores that order of explanation by starting from the ontology and then stating the role of \TheoryName{} inside the broader \TheoryProgram{}.

This paper adopts a strict claim boundary. It does not present a completed source derivation of the gravitational sector, and it does not claim that a selected research architecture is an established physical ontology. Instead, it advances four narrower claims. First, the selected source arena and its primitive debt can be stated exactly. Second, \(\PhiSrc\) can be retained as an unresolved slot without silently choosing its semantics. Third, \TheoryName{} provides the exact-GR benchmark used at the effective level. Fourth, historical interpretive language can be preserved as motivation without being laundered into canonical source structure or derivational success.

The sequence role of the present manuscript should therefore be read explicitly. After the overview and the short exact-closure anchor, this foundations paper opens the modular full statement by fixing the vocabulary that the later dynamics, consistency, relativistic-recovery, and flow-geometry papers will use. In that order, exact closure is primary and derivational continuation is secondary.

\section{Framework Statement}

\begin{definition}[\TheoryProgram{}]
\TheoryProgram{} is the umbrella framework that attempts to turn the \Aether{} / \Aflow{} ontology into a full theory program with explicit kinematics, dynamics, consistency conditions, and empirical interpretation.
\end{definition}

\begin{proposition}[Active theory status]
At the present stage of the repository, \TheoryName{} is the active theory developed in full manuscript form.
\end{proposition}

This proposition fixes the conceptual order of exposition in the main manuscript sequence. The active root is organized by exact closure rather than by an unfinished low-energy alternative proposal.

\section{Original Intent}

The enduring idea behind \TheoryProgram{} is not that gravity must immediately deviate from GR. It is that the relativistic world accessible to observers may be the effective manifestation of a deeper order. Earlier formulations expressed that intuition through the image of a three-dimensional active medium. The preserved insight in that language was that time is better understood as ordered change than as an independently existing container. The weakness of that language was that it encouraged a naive picture of an ordinary three-dimensional medium moving into an external emptiness.

The selected formulation replaces that mixed picture with a narrower source architecture. It assumes a smooth four-dimensional source arena and leaves the relation to observed space, expansion, and temporal order to later bridge work. The canonical map is therefore only
\[
\text{\Aether{}}
=
\text{continuum-first source arena } \Sub
\text{ (explicit primitive debt)},
\]
\[
\text{\Aflow{}}
=
\text{unresolved abstract slot } \PhiSrc.
\]
No arrow from this pair to observed space, physical time, matter, or exact-GR
geometry is adopted here. The older experiential-slice, S-time, and
deeper-motion phrases remain historical interpretation targets, not equations
or source definitions.

The scientific importance of this reordering is methodological. Ontology is no longer allowed to masquerade as completed dynamics. Instead, ontology fixes what the framework is trying to describe, while the later mathematical manuscripts must show how that description is closed, controlled, and related to established physics.

\section{Ontology of the \Aether{} and the \Aflow{}}

\begin{definition}[\Aether{}]
In the canonical mainline research architecture, the \Aether{} names a smooth
four-dimensional source manifold \(\Sub\). Its manifold character,
four-dimensionality, topology, and smooth structure are explicit primitive
debt. This definition does not identify \(\Sub\) with physical spacetime or
claim that the source arena is derived, unique, fundamental, or empirically
established.
\end{definition}

\begin{definition}[\Aflow{}]
At canonical source level, the \Aflow{} vocabulary is represented only by the
unresolved abstract slot \(\PhiSrc\). No parameter domain, group, semigroup,
monoid, invertibility condition, stochastic law, gauge meaning, generator, or
physical-time interpretation is adopted. It is not an ordinary vector wind in
observed three-dimensional space.
\end{definition}

These definitions carry an explicit rather than concealed burden. The selected
source package begins with
\[
\Sub,
\qquad
\dim \Sub = 4,
\]
together with the symbol \(\PhiSrc\) for a later source-order or evolution
specification. P5-T01 must type the surrounding source-theory slots, and
P5-T02 must resolve the semantics of \(\PhiSrc\). Until then, even writing a
map \(\Phi_\lambda:\Sub\to\Sub\) would add parameter and action structure that
the protected selection did not adopt.

This notation is intentionally modest. It does not define source dynamics,
degrees of freedom, a generator, observables, admissible variations, or a
reconstruction map. In particular, the \Aflow{} should not be reduced to a
naive preferred-frame vector, a target-side congruence, or a clock unless a
later authorized derivation supplies the relevant bridge.

\section{Why the \Aflow{} Is Not a Naive Vector Wind}

The flow language of the framework must be handled under strict control. If the \Aflow{} were treated as nothing more than an ordinary vector wind in observed space, the framework would risk importing precisely the wrong kind of low-energy structure: spurious preferred-frame signals, the wrong symmetry interpretation, and an overly literal picture of the ontology.

For that reason, the present manuscript uses flow language in a restricted
sense. At source level it names the unresolved slot \(\PhiSrc\), not completed
motion, a low-energy observable field, or a target-side congruence. Vectorial
or congruence-like representations may still be useful inside the exact-GR
interpretive dictionary, but no source-to-congruence map is adopted here.

This discipline is especially important in \TheoryName{}. The exact-closure theory retains the ontological role of the \Aflow{} while denying that it has already appeared as an extra established weak-field or radiative sector beyond GR.

\section{Observed Space as Local Experiential Slice}

The phrase \emph{observed space as a local experiential slice} is preserved as
historical interpretive motivation, not as an adopted source structure. The
canonical package supplies no embedding
\(\iota_s:\Sigma_s\hookrightarrow\Sub\), no source metric \(G_{AB}\), and no
induced spatial metric. Because \(\Sub\) is not identified with physical
spacetime, such an embedding cannot be inferred from the shared dimension or
smoothness alone.

A later bridge may attempt to construct observer-accessible spatial structure
from source data. It must type the source and target objects independently and
show why the map is lawful without importing the target atlas, metric, or
benchmark behavior as a premise. Until that burden is discharged,
\emph{experiential slice} is interpretation language only.

\section{S-Time as Experienced Order of Change}

The concept of S-time preserves an early interpretive aim of the program:
relating experienced temporal order to source organization without treating
time as a place-like extra corridor. It is not part of the selected primitive
package. In particular, the canonical source architecture supplies no
functional \(\Order[\Gamma,\Psi,\Sub,\PhiSrc]\), no monotonicity law, and no
identification of a source-order parameter with physical clock time.

The phrase \emph{experienced time} or \emph{S-time} may therefore be used only
as a name for an open reconstruction target. P5-T02 must first resolve the
source semantics of \(\PhiSrc\); a later bridge must separately construct
operational clocks, causal orientation, and calibration.

\section{Observed Expansion, Light, and Relativistic Phenomenology}

The framework does not deny the observational reality of expansion, clock behavior, redshift, or relativistic duration effects. It relocates their explanatory status. What is primary in the ontology is the deeper ordered substrate; what is primary in observation is the effective relativistic structure accessible to embedded observers.

The selected source architecture does not yet supply a lawful implication from
\((\Sub,\PhiSrc)\) to evolving slice geometry, light propagation, temporal
ordering, or observed expansion. Those are target-side phenomena governed in
\TheoryName{} by the adopted exact-GR benchmark. A later first-principles
recovery must construct the relevant source-to-target bridge without importing
the target metric or benchmark behavior as source evidence.

The present foundations paper does not claim that those phenomena have already been derived from substrate microphysics. Its narrower claim is that the ontology is compatible with their exact recovery and that, in \TheoryName{}, that recovery is secured by exact closure to GR.

\section{\TheoryName{} as the Exact-Closure Realization of \TheoryProgram{}}

\begin{definition}[\TheoryName{}]
\TheoryName{} is the exact-closure, GR-consistent benchmark of
\TheoryProgram{} in which the selected source research architecture and its
interpretive vocabulary remain distinct from the effective gravitational
dynamics, which are taken to be exactly those of general relativity with
ordinary matter coupling.
\end{definition}

This theory is exact closure rather than deformation. Its effective action is
\begin{equation}
S_{\mathrm{eff}}
=
\frac{c^3}{16\pi G}\int d^4x\,\sqrt{-g}\,R
+
S_{\mathrm{matter}}[g,\psi],
\label{eq:SA}
\end{equation}
with field equations
\begin{equation}
G_{\mu\nu} = \frac{8\pi G}{c^4} T_{\mu\nu}.
\label{eq:einstein}
\end{equation}
Null propagation and proper-time measurement are therefore governed by the standard relativistic relations
\begin{equation}
0 = g_{\mu\nu}\,dx^{\mu}dx^{\nu},
\label{eq:null}
\end{equation}
\begin{equation}
d\tau^2 = -\frac{1}{c^2} g_{\mu\nu}\,dx^{\mu}dx^{\nu}.
\label{eq:propertime}
\end{equation}

\begin{proposition}
The predictive content of \TheoryName{} is exactly that of GR. Accordingly, \TheoryName{} carries no independent low-energy non-GR observable signature.
\end{proposition}

\begin{remark}
This does not make \TheoryName{} scientifically empty. It makes its status precise. The theory is a disciplined interpretive completion within \TheoryProgram{} and a proposed deeper ontology consistent with GR, not a distinct low-energy alternative to it.
\end{remark}

\section{\TheoryName{} as Benchmark and Reversion Theory}

\begin{proposition}[Benchmark role]
Any future extension of \TheoryProgram{} must recover \TheoryName{} in the appropriate limit. \TheoryName{} therefore functions as the internal exact benchmark of the full framework.
\end{proposition}

\begin{proposition}[Automatic reversion]
If a proposed extension of the framework fails an essential consistency or viability gate, the scientific status of the framework contracts back to \TheoryName{} rather than collapsing into incoherence.
\end{proposition}

This benchmark structure is one of the central reasons to make \TheoryName{} the active manuscript line now. A disciplined theory program cannot allow an unfinished deviation sector to remain the narrative center once the exact benchmark theory is already available and the currently tested deviation theory has not earned promotion.

The correct short status statement is therefore the following: \TheoryName{} is the exact-closure theory developed in the present paper and is the active reference theory of the repository.

\section{Relation to General Relativity}

The relation between \TheoryName{} and GR is exact at the operational level and interpretive at the ontological level. Operationally, GR remains the governing relativistic theory of gravitation in this theory. Ontologically, \TheoryProgram{} adds the proposal that the relativistic metric description may be the effective manifestation of a deeper ordered substrate.

Two clarifications are essential. First, \TheoryName{} does not dispute any tested relativistic phenomenon. Clock behavior, null propagation, redshift, time dilation, and weak-field structure are all inherited from GR exactly through \eqref{eq:SA}--\eqref{eq:propertime}. Second, selection of the continuum-first source architecture does not license the introduction of an extra measured preferred-frame field in the effective theory. The source architecture and its interpretation remain an open derivational program, not an additional established low-energy sector.

Local special-relativistic kinematics are also retained through this exact closure. Since the effective metric sector is Einsteinian, the local inertial description continues to reproduce the standard special-relativistic content seen by observers. The full relation to SR should still be treated more explicitly in a later relativistic-recovery manuscript, but the foundations-level point is already clear: the deeper ontology is not offered as a contradiction of relativity, but as a possible deeper account of why an effective relativistic structure governs observation.

\section{Claim Boundary and Program Status}

The present paper separates five levels of statement that should not be blurred together.

\subsection{Ontology}

The canonical mainline research architecture takes a smooth
four-dimensional source manifold as explicit primitive debt and retains only
an unresolved abstract \(\PhiSrc\) slot. This is a tracked architecture
selection, not a physical-ontology truth claim. Experiential-slice and S-time
language remains interpretation and open reconstruction work.

\subsection{Exact closure}

The framework claims that \TheoryName{} provides an exact relativistic closure in which the operational content of GR is recovered without modification.

\subsection{Adopted dynamics}

The framework adopts the Einstein equations as the effective dynamics of gravity in \TheoryName{}. This is an adoption claim and should not be misdescribed as a completed substrate derivation.

\subsection{Open derivational work}

What remains unfinished is the recovery of the effective relativistic sector from explicit substrate variables, a closed substrate action, controlled matter coupling, mode analysis, and consistency proofs. Those tasks belong to the later dynamics and consistency manuscripts.

\subsection{Established in the present paper}

The following points are claimed as established at the present foundations level:
\begin{itemize}
    \item the protected selection and canonical integration of the continuum-first source research architecture;
    \item the four-dimensional smooth source arena as explicit primitive debt rather than a derived or physical spacetime;
    \item the unresolved status of \(\PhiSrc\) and the non-adoption of its possible semantics;
    \item the theory-level definition of \TheoryName{} as the exact-closure realization of \TheoryProgram{};
    \item the exact operational relation of \TheoryName{} to GR;
    \item the benchmark and reversion role of \TheoryName{} inside the total framework.
\end{itemize}

\subsection{Not yet established}

The following stronger statements are \emph{not} claimed here:
\begin{itemize}
    \item a completed substrate derivation of Einsteinian dynamics;
    \item a full microphysical definition of the substrate degrees of freedom;
    \item any adopted source-to-observer slice embedding or S-time functional;
    \item a domain, algebra, invertibility rule, stochastic law, gauge meaning, generator, or physical-time meaning for \(\PhiSrc\);
    \item a completed consistency analysis of a specific substrate action;
    \item a completed alternative low-energy sector beyond the present exact-closure line.
\end{itemize}

\section{Conclusion}

This foundations manuscript now provides the opening module of the fuller
public-facing statement of \TheoryProgram{}. Its central claims are limited
but definite. The canonical mainline uses a smooth four-dimensional source
manifold as explicit primitive debt and leaves \(\PhiSrc\) unresolved. It does
not identify that source arena with physical spacetime or adopt a slice map,
clock, metric, matter semantics, coupling, source law, or source dynamics.
Within that boundary, \TheoryName{} remains the exact-closure benchmark whose
effective gravitational dynamics are exactly those of GR.

The result is a stable starting point rather than a completed final theory. It
tells the reader what the \TheoryProgram{} is about, what \TheoryName{} is,
why the \Aflow{} should not be read as a naive vector wind, which source
assumptions are primitive, why \TheoryName{} is the active reference theory,
and what the framework presently adopts rather than derives. It therefore
secures the source-architecture and conceptual ground on which the
exact-closure benchmark can be stated without confusing selection, adoption,
interpretation, and derivation.

\clearpage
\begingroup
\renewcommand{\refname}{References}
\makeatletter
\renewcommand{\@biblabel}[1]{}
\makeatother
\begin{thebibliography}{99}
\small
\setlength{\itemsep}{0.25em}
\setlength{\parskip}{0pt}

\bibitem{Einstein1916}
Einstein, A. (1916). The foundation of the general theory of relativity. \emph{Annalen der Physik}, \emph{49}, 769--822.

\bibitem{ADM1962}
Arnowitt, R., Deser, S., \& Misner, C. W. (1962). The dynamics of general relativity. In L. Witten (Ed.), \emph{Gravitation: An introduction to current research} (pp. 227--265). Wiley. Reprinted in \emph{General Relativity and Gravitation}, \emph{40}, 1997--2027.

\bibitem{Raychaudhuri1955}
Raychaudhuri, A. (1955). Relativistic cosmology. I. \emph{Physical Review}, \emph{98}, 1123--1126.

\bibitem{EllisVanElst1999}
Ellis, G. F. R., \& van Elst, H. (1999). Cosmological models (Carg\`ese lectures 1998). In M. Lachi\`eze-Rey (Ed.), \emph{Theoretical and observational cosmology} (NATO Science Series C: Mathematical and Physical Sciences, Vol. 541, pp. 1--116). Kluwer. \href{https://arxiv.org/abs/gr-qc/9812046}{arXiv:gr-qc/9812046}

\bibitem{Riess1998}
Riess, A. G., et al. [Supernova Search Team]. (1998). Observational evidence from supernovae for an accelerating universe and a cosmological constant. \emph{Astronomical Journal}, \emph{116}, 1009--1038. \href{https://arxiv.org/abs/astro-ph/9805201}{arXiv:astro-ph/9805201}

\bibitem{Perlmutter1999}
Perlmutter, S., et al. [Supernova Cosmology Project]. (1999). Measurements of $\Omega$ and $\Lambda$ from 42 high-redshift supernovae. \emph{Astrophysical Journal}, \emph{517}, 565--586. \href{https://arxiv.org/abs/astro-ph/9812133}{arXiv:astro-ph/9812133}

\bibitem{Planck2018Parameters}
Aghanim, N., et al. [Planck Collaboration]. (2020). Planck 2018 results. VI. Cosmological parameters. \emph{Astronomy \& Astrophysics}, \emph{641}, A6. \href{https://arxiv.org/abs/1807.06209}{arXiv:1807.06209}

\bibitem{JacobsonMattingly2001}
Jacobson, T., \& Mattingly, D. (2001). Gravity with a dynamical preferred frame. \emph{Physical Review D}, \emph{64}, 024028.

\bibitem{JacobsonMattingly2004}
Jacobson, T., \& Mattingly, D. (2004). Einstein-aether waves. \emph{Physical Review D}, \emph{70}, 024003.

\bibitem{ElingJacobson2004}
Eling, C., \& Jacobson, T. (2004). Static post-Newtonian equivalence of general relativity and gravity with a dynamical preferred frame. \emph{Physical Review D}, \emph{69}, 064005.

\bibitem{Jacobson2010}
Jacobson, T. (2010). Extended Ho\v{r}ava gravity and Einstein-aether theory. \emph{Physical Review D}, \emph{81}, 101502(R).

\bibitem{Jacobson2014}
Jacobson, T. (2014). Undoing the twist: The Ho\v{r}ava limit of Einstein-aether theory. \emph{Physical Review D}, \emph{89}, 081501(R).

\bibitem{JacobsonSperanza2015}
Jacobson, T., \& Speranza, A. J. (2015). Variations on an aethereal theme. \emph{Physical Review D}, \emph{92}, 044030.

\bibitem{Horava2009}
Ho\v{r}ava, P. (2009). Quantum gravity at a Lifshitz point. \emph{Physical Review D}, \emph{79}, 084008.

\bibitem{Jacobson1995}
Jacobson, T. (1995). Thermodynamics of spacetime: The Einstein equation of state. \emph{Physical Review Letters}, \emph{75}, 1260--1263.

\bibitem{Visser2002}
Visser, M. (2002). Sakharov's induced gravity: A modern perspective. \emph{Modern Physics Letters A}, \emph{17}, 977--992. \href{https://arxiv.org/abs/gr-qc/0204062}{arXiv:gr-qc/0204062}

\bibitem{HojmanKucharTeitelboim1976}
Hojman, S. A., Kucha\v{r}, K., \& Teitelboim, C. (1976). Geometrodynamics regained. \emph{Annals of Physics}, \emph{96}, 88--135.

\bibitem{Deser2004}
Deser, S. (1970). Self-interaction and gauge invariance. \emph{General Relativity and Gravitation}, \emph{1}, 9--18. Reprinted with 2004 note at \href{https://arxiv.org/abs/gr-qc/0411023}{arXiv:gr-qc/0411023}

\bibitem{Lovelock1971}
Lovelock, D. (1971). The Einstein tensor and its generalizations. \emph{Journal of Mathematical Physics}, \emph{12}, 498--501.

\bibitem{BarceloLiberatiVisser2011}
Barcelo, C., Liberati, S., \& Visser, M. (2011). Analogue gravity. \emph{Living Reviews in Relativity}, \emph{14}, 3.

\bibitem{Sakharov1967}
Sakharov, A. D. (1968). Vacuum quantum fluctuations in curved space and the theory of gravitation. \emph{Soviet Physics Doklady}, \emph{12}, 1040--1041. Translated from \emph{Doklady Akademii Nauk SSSR}, \emph{177}, 70--71.

\end{thebibliography}
\endgroup


\end{document}
